% \usepackage{color}
% \definecolor{light-gray}{gray}{0.95}
% \definecolor{dark-green}{rgb}{0,0.25,0}
\usepackage{fontspec,xunicode,xltxtra}
\setmainfont[Mapping=tex-text]{Georgia}
% \setmainfont[Ligatures=Required,Mapping=tex-text]{Palatino Linotype}
% \setsansfont[Ligatures=Required,Mapping=tex-text]{Myriad Pro}
\setsansfont[Mapping=tex-text]{Arial}
\setmonofont[Mapping=tex-text]{Courier New}

\usepackage{xeCJK}
\setCJKmainfont[ItalicFont={Adobe Kaiti Std}]{Adobe Song Std}
% \setCJKmainfont[ItalicFont={Adobe Kaiti Std}]{Adobe Kaiti Std}
\setCJKsansfont{Adobe Heiti Std}
% \setCJKsansfont{Microsoft YaHei}
\setCJKmonofont{Adobe Heiti Std}
\punctstyle{banjiao}


\usepackage{indentfirst}


\usepackage{xcolor}
\definecolor{light-gray}{gray}{0.95}

%生成PDF的链接
\usepackage[dvipdfmx]{hyperref}
\hypersetup{
    bookmarks=true,         % show bookmarks bar?
    bookmarksopen=true,
    pdfpagemode=UseNone,    % options: UseNode, UseThumbs, UseOutlines, FullScreen
    pdfstartview=FitH,
    pdfborder=1,
    pdfhighlight=/P,
    pdfauthor={wuxch},
    unicode=true,           % non-Latin characters in Acrobat’s bookmarks
    colorlinks,             % false: boxed links; true: colored links
    linkcolor=blue,         % color of internal links
    citecolor=blue,        % color of links to bibliography
    filecolor=magenta,      % color of file links
    urlcolor=cyan           % color of external links
}
% 使用目录
\usepackage{makeidx}
% 使用图形
\usepackage{graphicx}


\usepackage{titletoc}

\usepackage[sf]{titlesec}

\dottedcontents{section}[1.5em]{}{2.3em}{1pc}


\titleformat*{\section}{\sffamily\bfseries\Large}

% % 页眉页脚的设定
% \usepackage{fancyhdr}
% \pagestyle{fancy}
% \lhead{\includegraphics[height=0.2in]{zte_logo}}

% \usepackage{footmisc}
% \usepackage{pfnote}
% \usepackage{fnpos}

% % 嵌入的代码显示
% \usepackage{listings}
% \lstset{language=C++, breaklines, extendedchars=false}
% \lstset{keywordstyle=\color{blue},
%         commentstyle=\color{SeaGreen},
%         stringstyle=\ttfamily,
%         showstringspaces=false,
%         backgroundcolor=\color{light-gray}}
% \lstset{basicstyle=\small}

\usepackage{booktabs}

\usepackage{tipa}
\usepackage{phonetic}
\usepackage[gen]{eurosym}

% \usepackage{CJK}
\usepackage{makeidx}
% \usepackage{glossary}

\usepackage{fancyhdr}

\usepackage{lastpage}
\cfoot{\textsf{\thepage\ of \pageref*{LastPage}}}

\pagestyle{fancy}
\fancyhead{}
\renewcommand{\headrulewidth}{0pt}
\pagenumbering{arabic}

\usepackage{graphicx}
\usepackage{picinpar}

\setlength{\parindent}{0in}
% \setlength{\parskip}{2.5ex plus4mm minus3mm}
\setlength{\parskip}{3ex plus4mm minus3mm}
% \setlength{\parskip}{.3\baselineskip}

\setcounter{section}{0}

% bibtex
\usepackage[super]{natbib}
% \usepackage[globalcitecopy,labelstoglobalaux]{bibunits}
% \makeatletter
% \renewcommand{\@bibunitname}{\jobname.\the\@bibunitauxcnt}
% \makeatother

% \renewcommand{\refname}{缩略语}
% \renewcommand{\refname}{}

\makeatletter                   % Reference list option change
\renewcommand\@biblabel[1]{#1:  } % from [1] to 1: 
\makeatother                    %
% 文章内的数字上标字体
\renewcommand{\citenumfont}[1]{\textsf{#1}}
% 文章后的数字字体
\renewcommand{\bibnumfmt}[1]{\textsf{#1}:}

\renewcommand{\refname}{Vocabulary}

\newcommand{\phonetic}[1]{{\fontspec[Colour=CC00CC]{Microsoft Sans Serif} #1}}

\newcommand{\vocabulary}{\pagebreak \bibliographystyle{d:/wuxch/emacs_home/vocabulary/wuxchvcb}
  \bibliography{d:/wuxch/emacs_home/vocabulary/vocabulary}}

% 首字下沉
\definecolor{Goldenrod}{rgb}{0.8549,0.6471,0.1255}
\usepackage{lettrine}
\setlength{\DefaultFindent}{0.5em}
\setlength{\DefaultNindent}{0em}
\renewcommand{\LettrineFontHook}{\color{Goldenrod}}

\usepackage{wuxchcal}

\newcommand{\mycalendar}[2]{\marginpar{\raisebox{-3.4ex}{\wuxchcal{1}{#1}{0.45}{#2}{1.3}}}}

% \newcommand{\mycalendar}[2]{\marginpar[aaa]{aaa}}

%%% Local Variables: 
%%% mode: latex
%%% TeX-master: "nyt16"
%%% End: 
